\label{sec:legal}

\subsection{Functional Economic Areas and Communities of Interest}

The concept of the functional economic area---a geographic region whose
residents share a labor market, a retail trade area, and a common daily
economic sphere---has been recognized in redistricting law as a legitimate
basis for community of interest claims.

The National Conference of State Legislatures' 2021 survey of redistricting
criteria found that 35 states include communities of interest as a criterion,
and that ``economic interests'' and ``shared economic community'' are among
the most frequently cited community subtypes in commission guidelines and
judicial interpretations \citep{ncsl2021criteria}.
A commuting shed is a precise quantification of the ``shared economic
community'' concept: the residents of a commuting shed literally share
their economic daily life---they compete for the same jobs, drive the same
routes, patronize the same establishments, and depend on the same employers.

\paragraph{California commuting communities.}
California's Citizens Redistricting Commission criteria cite communities
defined by ``economic and social interests'' \citep{caprop11}.
In its 2011 and 2021 proceedings, the Commission accepted testimony from
commuting communities:
residents of the Inland Empire testified that their communities shared
commuting patterns into the Los Angeles basin, creating economic
interdependencies not reflected in county boundaries.
Silicon Valley workers testified that the commuting shed from San Jose
through Mountain View to Menlo Park defined an economic community distinct
from the agricultural Santa Clara Valley south of Morgan Hill.
The M.4 algorithm operationalizes these commuting community arguments:
tracts with overlapping commuting destinations receive high similarity
and heavy edges; tracts commuting to different destinations receive low
similarity and light edges.

\paragraph{Colorado trade area communities.}
Colorado's redistricting criteria explicitly recognize ``trade area'' community
membership---communities that share a retail trade zone or labor market area
\citep{coart5s44}.
A commuting shed is the labor market equivalent of a trade area: it identifies
which residential communities are economically integrated with which employment
centers.
The Colorado Front Range provides a clear example: communities in the Boulder
valley commute to Boulder and Denver; communities in the Fort Collins corridor
commute to Fort Collins and the I-25 tech corridor.
These are distinct commuting sheds that correspond to distinct communities
of interest in redistricting proceedings.

\subsection{Airport and Port Communities}

Airport and port communities provide some of the clearest examples of
commuting-shed community identification in redistricting proceedings.

\paragraph{Airport communities.}
The area surrounding a major commercial airport is a distinctive commuting
community.
Residents of tracts near Chicago O'Hare International Airport, Houston
George Bush Intercontinental Airport, or Los Angeles International Airport
share an employment base (aviation industry, cargo logistics, hotel and
hospitality services) and commuting pattern.
The LODES OD data captures this concentration: workers in airport-adjacent
tracts send flows to the airport employment blocks, creating a tight
destination cluster.
Under M.4, airport tracts share high weighted Jaccard similarity with each
other and low similarity with non-airport tracts commuting to downtown office
employment.

This commuting shed community has been recognized in redistricting proceedings.
In the 2001 California redistricting, residents near LAX testified before
the Special Master that their community was defined by the airport economy---not
by their city of residence (Inglewood, El Segundo, and sections of Los Angeles)
or by the county boundary (all three are in Los Angeles County).
The court-appointed Special Master's final map kept the airport industrial and
residential community together, in a configuration consistent with M.4's
prediction of high within-shed similarity.

\paragraph{Port communities.}
Port communities around major shipping facilities have a similarly distinctive
commuting character.
Residents near the Port of Long Beach/Los Angeles, the Port of Houston-Galveston,
and the Port of Baltimore-Newport News all exhibit commuting patterns directed
primarily to the port and its logistics hinterland rather than to central
business districts.
Under M.4, these tracts cluster together based on their shared destinations
(port terminals, distribution centers, rail yards) and are appropriately
separated from adjacent residential tracts commuting to downtown employment.

The legal recognition of port communities as a distinct economic community of
interest has been affirmed in multiple redistricting proceedings.
The Maryland redistricting commission in 2011 accepted testimony from port
and waterfront industrial community advocates who argued that Baltimore's
waterfront industrial tracts formed a distinct community.
The final plan maintained the port community's cohesion in a way consistent
with the M.4 prediction.

\subsection{Partisan Neutrality of Commuting Data}

LODES OD data contains no electoral, racial, or partisan information.
The commute flow records are derived from UI wage records matched to
employer geocodes; they record only where workers live and where they work.
No information about the political preferences, racial composition, or
electoral behavior of the workers is present in the data.

Any partisan effect of commuting-shed weights---such as the polycentric
separation of Raleigh and Charlotte metros in the NC prediction---is a
consequence of the geographic correlation between commuting patterns and
partisan sorting.
Urban employment centers that attract commuters from diverse residential areas
are a feature of the population landscape, not a product of the algorithm.

Following \textit{Rucho v.\ Common Cause} \citeyearpar{rucho2019}, the use
of a non-partisan, federally administered dataset to operationalize a
legitimate redistricting criterion (communities of interest) satisfies the
constitutional standard for non-partisan redistricting even if the resulting
plan has partisan effects.

\subsection{Expert Witness Deployment Template}

Redistricting practitioners deploying M.4 commuting shed weights should
structure expert testimony around the following elements:

\paragraph{1. Data provenance declaration.}
``I used the U.S.\ Census Bureau's LODES Origin-Destination Employment
Statistics, version~8, for census year 2020
(\url{https://lehd.ces.census.gov/data/}).
This is a federal administrative dataset derived from state Unemployment
Insurance wage records matched to residential and workplace geocodes.
It records, for each home census block, how many workers residing in
that block are employed in each work census block.
No voter registration data, election results, racial composition data,
or any other socioeconomic information about workers is present in the
dataset or used in the computation.''

\paragraph{2. Commuting community identification.}
``I computed a commuting destination distribution for each census tract:
a record of how many workers residing in each tract commute to each
employment destination.
Two census tracts are in the same commuting shed if they share a substantial
fraction of their commuting destinations.
This overlap is measured by the weighted Jaccard similarity coefficient:
the flow-weighted intersection of commuting destinations divided by the
flow-weighted union.
District boundaries are drawn preferentially at commuting shed boundaries:
tracts with overlapping destinations are kept together; tracts commuting
to different economic areas are separated.''

\paragraph{3. Specific community identification.}
``In [state], the following commuting communities are preserved by the
proposed plan: [specific examples, e.g., `the Research Triangle Park
commuting shed, comprising tracts in Wake, Durham, Orange, and Chatham
counties whose workers commute to the RTP technology corridor, is preserved
as a cohesive economic community within two contiguous congressional districts.
The Charlotte metro commuting shed, comprising Mecklenburg and Union county
tracts commuting to the Uptown Charlotte and I-77/I-85 employment corridors,
is preserved separately.
No district combines Research Triangle and Charlotte commuting shed tracts,
reflecting the 160-mile functional economic separation between these areas'].''

\paragraph{4. Criterion alignment.}
``The [state]'s redistricting criteria require preservation of communities
defined by economic and shared-interest considerations.
The commuting shed similarity weights directly implement this criterion:
communities of workers who share a functional economic area are kept
together by the higher edge weights between their residential tracts.''
